\DocSection{Corrected Inference Method}
\label{sec:corrected-method}

\DocSubsection{Direct stochastic variational objective}

Let \(q_\theta\) be a tractable distribution over complete deterministic
policies. The corrected method maximizes the evidence lower bound implied by
Equation~\eqref{eq:gibbs-variational}:
\begin{equation}
  \mathcal{L}(\theta)
  = \beta\E_{\policy\sim q_\theta}
      [\ExpectedReturn_H(\policy)]
    - \KL(q_\theta\|\mu).
  \label{eq:correct-elbo}
\end{equation}
Sampling \(\policy\sim q_\theta\) and then
\(\tau\sim p(\cdot\mid\policy)\) gives an unbiased estimator of the
return term because Equation~\eqref{eq:correct-elbo} is linear in
\(\ExpectedReturn_H\).

For a score-function estimator, a valid single-sample gradient contribution is
\begin{equation}
  \left(
    \beta\Return_H(\tau)
    + \log\mu(\policy)
    - \log q_\theta(\policy)
    - b
  \right)
  \nabla_\theta\log q_\theta(\policy),
  \label{eq:score-gradient}
\end{equation}
where \(b\) is an action-independent control variate. Multiple independent
rollouts of the same policy may be averaged before applying
Equation~\eqref{eq:score-gradient}; this reduces variance without changing the
target.

\DocSubsection{Policy representation and execution}

For a finite tabular policy, actions may be sampled lazily and memoized on the
first visit to each state. The log proposal probability contains each sampled
state assignment exactly once. Particle resampling, when used as an auxiliary
search mechanism, must copy the state, policy memo, visit counters, and any
other ancestry-dependent state together.

The primary execution rule samples one policy-level object at the beginning of
an episode and preserves it until termination. A separately named
\emph{marginal-action} ablation samples from estimated statewise marginals on
every visit. Comparing these two rules measures the value of policy-level
correlation instead of assuming it away.

\DocSubsection{Role of common random numbers}

Shared random streams may compare candidate policies under matched transition
noise. They are treated as a variance-reduction design, not as part of the
probabilistic target. A correct estimator should converge to the same result
under shared and independent streams; finite-sample efficiency may differ.

\DocSubsection{Replicated sample-average policy SMC}

For a fixed bank of \(K\) simulator tapes
\(\Xi_K=(\xi_1,\ldots,\xi_K)\), define
\begin{equation}
  \widehat J_K(\policy;\Xi_K)
  = \frac{1}{K}\sum_{k=1}^{K}\Return_H(\policy,\xi_k).
  \label{eq:saa-return}
\end{equation}
The conditional sample-average target is
\begin{equation}
  P_{\beta,K}(d\policy\mid\Xi_K)
  \propto \exp\!\left(\beta\widehat J_K(\policy;\Xi_K)\right)
  \mu(d\policy).
  \label{eq:saa-target}
\end{equation}

\DocSubsection{Learned full-support proposal}

Let \(m\) denote the assignments already stored in a particle's policy memo,
and let \(\mathcal A_\mu(s,m)\) be the prior-supported legal actions. A
stochastic policy \(q_{\mathrm{PPO},\phi}\), trained with PPO, is used as a
guide rather than as the inference target. With \(u_\mu(\cdot\mid s,m)\)
uniform on \(\mathcal A_\mu(s,m)\), the actual proposal is
\begin{equation}
  q_{\epsilon,\phi}(a\mid s,m)
  =
  (1-\epsilon)q_{\mathrm{PPO},\phi}(a\mid s,m)
  + \epsilon u_\mu(a\mid s,m),
  \qquad 0<\epsilon\leq 1.
  \label{eq:guided-proposal}
\end{equation}
Both components use the same legal-action mask. Hence every action in the
target support has proposal probability at least
\(\epsilon/|\mathcal A_\mu(s,m)|\). The mixture probability in
Equation~\eqref{eq:guided-proposal}, rather than the unmixed PPO probability,
is the denominator in every importance ratio.

Each particle contains one policy memo and \(K\) environment replicas. Let
\(U_t\) be the deterministically ordered set of distinct unassigned states
first exposed by those replicas at simulation step \(t\). One action is drawn
and memoized for each \(s\in U_t\). The exact incremental log-weight is
\begin{equation}
  \Delta\log w_t
  =
  \frac{\beta}{K}\sum_{k=1}^{K}r_{t,k}
  +
  \sum_{s\in U_t}
  \left[
    \log\mu(a_s\mid s,m_{<s})
    -\log q_{\epsilon,\phi}(a_s\mid s,m_{<s})
  \right].
  \label{eq:guided-weight}
\end{equation}
Thus the mean reward is formed before exponentiation, and each prior/proposal
correction is applied once per newly instantiated policy assignment, not once
per replica or visit. Revisits reuse the memoized action. Resampling copies all
replica states, the complete policy memo, and proposal context. No transition
ratio is required because the conditional target and proposal use the same
simulator dynamics.

The guide is frozen within each SMC sweep, and the parameters, action mask,
and value of \(\epsilon\) used to generate every proposal are recorded.
Proposal adaptation between sweeps, if studied, uses a separate frozen
snapshot for each sweep.

For finite \(K\), Equation~\eqref{eq:saa-target} is an explicit approximation,
not an unbiased representation of Equation~\eqref{eq:policy-target}. Under
bounded returns and a finite policy class it converges uniformly as \(K\)
increases. Experiments therefore report convergence in \(K\) separately from
convergence in the number of policy particles.

\DocSubsection{Proposal accounting and isolation}

Training \(q_{\mathrm{PPO},\phi}\) is part of the guided method. Matched-budget
comparisons count all simulator interactions, optimizer updates, wall-clock
time, accelerator time, and peak memory used to train the proposal as well as
the later SMC cost. Any amortized accounting is reported only as a separately
specified multi-instance setting.

The required isolation study includes uniform-proposal SMC
(\(\epsilon=1\)), corrected guided SMC using
Equations~\eqref{eq:guided-proposal}--\eqref{eq:guided-weight}, and a
proposal-alone condition. The proposal-alone condition samples the same lazy,
memoized deterministic policy from \(q_{\epsilon,\phi}\) but performs no
importance weighting or resampling. It therefore isolates the contribution of
SMC from that of the learned guide while preserving policy-commitment
semantics. These conditions define the evaluation protocol; they do not state
an empirical outcome.
